\section{Implementation}
\label{sec:implementation}

ALS is implemented as a standalone Rust crate (\texttt{audit-trail}) using
the Tokio async runtime. Table~\ref{tab:deps} lists the primary crate
dependencies. The four event sources each integrate an \texttt{als} module
that exposes \texttt{ALSClient} and \texttt{AlsQueue}; this module is
identical across all sources, keeping integration non-intrusive.

\begin{table}[htbp]
    \caption{Key ALS Dependencies}
    \label{tab:deps}
    \begin{center}
        \begin{tabular}{l l p{4cm}}
            \toprule
            \textbf{Crate} & \textbf{Version} & \textbf{Role} \\
            \midrule
            \texttt{axum}       & 0.8.4  & HTTP server for event ingestion \\
            \texttt{tokio}      & 1.44.2 & Async runtime \\
            \texttt{aes-gcm}    & 0.10   & AES-256-GCM encryption \\
            \texttt{p256}       & 0.13   & ECDH P-256 for ECIES key wrapping \\
            \texttt{iota-sdk}   & 1.x    & IOTA transaction construction \\
            \texttt{sha2}       & 0.10   & SHA-256 for hash-chaining and KDF \\
            \texttt{uuid}       & 1.x    & UUID v7 record identifiers \\
            \texttt{reqwest}    & 0.12   & HTTP client for IPFS upload \\
            \bottomrule
        \end{tabular}
    \end{center}
\end{table}

\subsection{Event Source Integration}

Integration into each event source is limited to adding an \texttt{als}
module call after each auditable operation. The call is non-blocking: it
enqueues an \texttt{AuditEvent} into the local \texttt{AlsQueue} and returns
immediately. The retry worker inside \texttt{ALSClient} handles delivery
failures transparently, retrying with exponential backoff up to a configured
limit before dropping the event and emitting a warning log.

\subsection{Move Smart Contract}

The on-chain component is a Move module \texttt{als::audit\_log} deployed
on the IOTA smart contract. Its entry function is:

\begin{verbatim}
public entry fun commit_rotation(
    cap: &AuditLogCap,
    cid: vector<u8>,
    file_hash: vector<u8>,
    record_count: u64,
    period_start: u64,
    period_end: u64,
    ctx: &mut TxContext
)
\end{verbatim}

The function creates a \texttt{RotationRecord} struct and immediately freezes
it via \texttt{transfer::freeze\_object}. A frozen Move object becomes
permanently read-only: no entry function---including one signed by the
\texttt{AuditLogCap} holder---can mutate it after finalization.
\texttt{AuditLogCap} is a one-time capability object held by ALS, preventing
unauthorized parties from publishing false rotation records.

\subsection{Implementation Scope}

Six implementation units (I-ALS-01 through I-ALS-06) map directly to the
design components described in Section~\ref{sec:design}: type schema
(I-ALS-01), event encapsulation (I-ALS-02), Audit Receiver (I-ALS-03),
Audit Logger with hash-chaining (I-ALS-04), rotation and archival worker
(I-ALS-05), and event source integration module (I-ALS-06). The
implementation was carried out in a prototype environment running Ubuntu
24.04 (WSL2) for ALS and client components, with a VPS hosting IPFS Cluster
and Gas Station, and IOTA Localnet for the smart contract.